\begin{figure}[H]
	\centering
	\includegraphics[width=0.25\textwidth]{./images/mnist_digit.png}
	\caption{Sample MNIST Digit (5)}
	\label{fig:example}
\end{figure}


\section{Data Loading}
The dataset we will be using to train a multi-layer neural network is the
MNIST dataset. The {MNIST}\footnote{\textit{Modified} National Institute of
	Standards and Technology} dataset is a collection of grayscale handwritten
digits commonly used for training various image processing systems. The dataset
contains 60,000 training images and 10,000 testing images each normalized to $28
	\times 28$ pixels.

\subsection{Data Splitting}
The MNIST dataset is split using an 80-20 ratio. Since the MNIST dataset
contains 60,000 training images and 10,000 testing images, we expect to have
48,000 images and target labels in the training set and 12,000 images and
target labels in the validation set. Optionally, the dataset is randomly
shuffled and seeded for reproducability.

\subsection{Normalization}
The training and validation datasets are normalized by z-score.

$$ \vec{z} = \frac{\vec{x} - \mu}{\sigma} $$

Here, $\vec{x}$ is the flattened ($ 28 \times 28 = 784$) vector, $\mu$ is
the mean pixel value, and $\sigma$ is the standard deviation over the single
image. The resulting vector $\vec{z}$ (also size $784$) contains the z-score
normalization of the image where each value will range between -1 and 1 (unit
variance) and have mean 0.

\subsection{Mean and Standard Deviation}
Consider an image of the digit $5$\cref{fig:example}. Its vector representation is

\begin{equation*}
	\vec{x}^{(n)}  =
	\begin{bmatrix}
		x_1    \\
		x_2    \\
		\vdots \\
		x_n
	\end{bmatrix}
\end{equation*}

where each of the values in $\vec{x}^{(784)}$ is some value between 0 for black
and 255 for white. Its mean ($\mu$) and standard deviation ($\sigma$) can be
obtained by

\begin{equation*}
	\begin{aligned}
		\mu    & = \frac{x_1 + x_2 + \ldots + x_n}{n}            & = 0.13714225924744897 \\
		\sigma & = \sqrt{\frac{\sum_{i=1}^{n} (x_i - \mu)^2}{n}} & = 0.31112823799846606
	\end{aligned}
\end{equation*}

These values can then be used to obtain $\vec{z}^{(n)}$, the z-score
normalization of $\vec{x}^{(n)}$.
